% =============================================================================
%  Section 2 — MCP-Assisted Evaluation (ByteBell Knowledge Graph)
%  Drop into the full report between Section 1 (Comparison) and Section 3 (Raw)
% =============================================================================

\section{MCP-Assisted Evaluation (ByteBell Knowledge Graph)}\label{sec:mcp}

\subsection{Methodology}

Each of the 22 \texttt{astropy/astropy} tasks from the
\texttt{princeton-nlp/SWE-bench\_Verified} dataset was presented to
Claude Code (model: \texttt{claude-opus-4-6}) with access restricted
\emph{exclusively} to ByteBell MCP tools.
The model could not read local files, run \texttt{git} commands, call
any remote REST/GraphQL API, or use web search.
All codebase understanding had to be derived through the ByteBell
knowledge graph --- a pre-indexed, semantically tagged representation of
the repository.

\paragraph{Prompt structure.}
Every task was issued under a minimal prompt with two components:

\begin{enumerate}
  \item \textbf{Hard constraint:} ``ONLY use ByteBell MCP tools --- no
        local file reads, no git APIs, no GitNexus skills.''
  \item \textbf{Answer format:} ``Just write your answer, no format
        needed.  \texttt{\{\ answer:\ "answer"\ \}}\ --- that's it.''
\end{enumerate}

\noindent
The SWE-Bench question text was appended directly after these rules.
Claude Code then explored the codebase through MCP tool calls
(\texttt{smart\_search}, \texttt{keyword\_lookup},
\texttt{retrieve\_file}, \texttt{graph\_traverse}, etc.) and wrote
its answer to \texttt{answer.json}.
No explicit instruction to produce a diff was given; the model chose
its own output format (descriptions, code blocks, or unified diffs)
per task.

\paragraph{Judging.}
Answers were graded against the ground-truth patch using the same
three-dimensional rubric as the raw runs (Table~\ref{tab:rubric-mcp}),
scored by a separate Claude Opus~4.6 judge session.

\begin{table}[h]
\centering
\caption{Scoring rubric --- 8-point scale (identical to raw runs).}\label{tab:rubric-mcp}
\begin{tabular}{lcp{8cm}}
\toprule
\textbf{Dimension} & \textbf{Max} & \textbf{Criterion} \\
\midrule
Root cause identification & 3 & Correctly diagnoses \emph{why} the bug exists. \\
Correct file(s)           & 2 & Identifies the right file(s) to change. \\
Correct patch / code change & 3 & Proposed change matches ground truth or is functionally equivalent. \\
\bottomrule
\end{tabular}
\end{table}

\noindent
Grade tiers: \textbf{Exact}~(8/8), \textbf{Near-perfect}~(7/8),
\textbf{Partial}~(5--6/8), \textbf{Fail}~($\le$4/8).

% ---------------------------------------------------------------------------
\subsection{Results}

\subsubsection{Aggregate Performance}

\begin{table}[h]
\centering
\caption{Dimension-level scores across all 22 MCP tasks.}\label{tab:dim-mcp}
\begin{tabular}{lrrr}
\toprule
\textbf{Dimension} & \textbf{Score} & \textbf{Max} & \textbf{\%} \\
\midrule
Root cause identification   & 64 & 66 & 97.0\% \\
Correct file(s)             & 44 & 44 & 100\%  \\
Correct patch / code change & 55 & 66 & 83.3\% \\
\midrule
\textbf{Overall}            & \textbf{163} & \textbf{176} & \textbf{92.6\%} \\
\bottomrule
\end{tabular}
\end{table}

\begin{table}[h]
\centering
\caption{Grade distribution (MCP).}\label{tab:grade-mcp}
\begin{tabular}{lcl}
\toprule
\textbf{Grade} & \textbf{Count} & \textbf{Tasks} \\
\midrule
Exact (8/8)        & 13 & 1--4, 7, 9, 12--14, 16--18, 21 \\
Near-perfect (7/8) &  6 & 5, 8, 10, 11, 20, 22 \\
Partial (5--6/8)   &  3 & 6, 15, 19 \\
Fail ($\le$4/8)    &  0 & --- \\
\bottomrule
\end{tabular}
\end{table}

\subsubsection{Per-Task Scores, Time \& Cost}

\begin{table}[h]
\centering
\small
\caption{Per-task MCP results --- score, wall-clock time, and API cost.}\label{tab:per-task-mcp}
\setlength{\tabcolsep}{4pt}
\begin{tabular}{clccccrr}
\toprule
\textbf{\#} & \textbf{Instance ID} & \textbf{RC} & \textbf{Files} & \textbf{Patch} & \textbf{Score} & \textbf{Time (s)} & \textbf{Cost (\$)} \\
\midrule
 1 & \texttt{astropy-12907} & 3 & 2 & 3 & 8/8 &    43 & 0.297 \\
 2 & \texttt{astropy-13033} & 3 & 2 & 3 & 8/8 &    66 & 0.239 \\
 3 & \texttt{astropy-13236} & 3 & 2 & 3 & 8/8 &   131 & 0.598 \\
 4 & \texttt{astropy-13398} & 3 & 2 & 3 & 8/8 &   281 & 0.697 \\
 5 & \texttt{astropy-13453} & 3 & 2 & 2 & 7/8 &    71 & 0.341 \\
 6 & \texttt{astropy-13579} & 3 & 2 & 1 & 6/8 &   472 & 0.501 \\
 7 & \texttt{astropy-13977} & 3 & 2 & 3 & 8/8 &   392 & 0.475 \\
 8 & \texttt{astropy-14096} & 3 & 2 & 2 & 7/8 &   134 & 0.440 \\
 9 & \texttt{astropy-14182} & 3 & 2 & 3 & 8/8 &    82 & 0.475 \\
10 & \texttt{astropy-14309} & 3 & 2 & 2 & 7/8 &    52 & 0.276 \\
11 & \texttt{astropy-14365} & 3 & 2 & 2 & 7/8 &    96 & 0.445 \\
12 & \texttt{astropy-14369} & 3 & 2 & 3 & 8/8 &   845 & 1.716 \\
13 & \texttt{astropy-14508} & 3 & 2 & 3 & 8/8 &   413 & 0.326 \\
14 & \texttt{astropy-14539} & 3 & 2 & 3 & 8/8 &   416 & 0.331 \\
15 & \texttt{astropy-14598} & 2 & 2 & 1 & 5/8 &   547 & 1.474 \\
16 & \texttt{astropy-14995} & 3 & 2 & 3 & 8/8 &    55 & 0.312 \\
17 & \texttt{astropy-7166}  & 3 & 2 & 3 & 8/8 &   145 & 0.604 \\
18 & \texttt{astropy-7336}  & 3 & 2 & 3 & 8/8 &    83 & 0.361 \\
19 & \texttt{astropy-7606}  & 2 & 2 & 2 & 6/8 &    94 & 0.352 \\
20 & \texttt{astropy-7671}  & 3 & 2 & 2 & 7/8 &   785 & 0.492 \\
21 & \texttt{astropy-8707}  & 3 & 2 & 3 & 8/8 &    84 & 0.394 \\
22 & \texttt{astropy-8872}  & 3 & 2 & 2 & 7/8 &   187 & 0.387 \\
\midrule
   & \textbf{Total}  & \textbf{64/66} & \textbf{44/44} & \textbf{55/66} & \textbf{163/176} & \textbf{5{,}474} & \textbf{11.53} \\
   & \textbf{Average} &       &       &       & \textbf{7.4/8}   & \textbf{249} & \textbf{0.52} \\
\bottomrule
\end{tabular}
\end{table}

% ---------------------------------------------------------------------------
\subsubsection{Token \& Cost Breakdown (Aggregated)}

\begin{table}[h]
\centering
\small
\caption{Aggregated token usage and cost across all 22 MCP tasks.}\label{tab:tokens-agg-mcp}
\setlength{\tabcolsep}{4pt}
\begin{tabular}{lrrrrrrr}
\toprule
& \textbf{Input} & \textbf{Cache Write} & \textbf{Cache Read} & \textbf{Output}
& \textbf{Eff Weighted} & \textbf{Eff Input} & \textbf{Cost (\$)} \\
\midrule
Total   & 50{,}518 & 704{,}002 & 7{,}874{,}141 & 126{,}060 & 1{,}717{,}935 & 8{,}628{,}661 & 11.53 \\
Average & 2{,}296  & 32{,}000  & 357{,}915     &   5{,}730 &    78{,}088   &   392{,}212   &  0.52 \\
\bottomrule
\end{tabular}
\\[4pt]
{\footnotesize
Effective Weighted Input = Input + $1.25 \times$ Cache Write + $0.1 \times$ Cache Read.\quad
Effective Input = Input + Cache Write + Cache Read.}
\end{table}

% ---------------------------------------------------------------------------
\subsubsection{Per-Model Token Usage}

As with the raw runs, Claude Code dispatches requests to
\texttt{claude-opus-4-6} and \texttt{claude-haiku-4-5}.
In the MCP workflow the split is even more extreme
(Table~\ref{tab:model-split-mcp}).

\begin{table}[h]
\centering
\caption{Token and cost split between Opus and Haiku (MCP).}\label{tab:model-split-mcp}
\begin{tabular}{lrrrr}
\toprule
\textbf{Metric} & \textbf{Opus} & \textbf{Haiku} & \textbf{Total} & \textbf{Opus \%} \\
\midrule
Input tokens       &        438 &    50{,}080 &     50{,}518 &   0.9\% \\
Cache Write tokens &  704{,}002 &          0 &    704{,}002 & 100\%   \\
Cache Read tokens  & 7{,}874{,}141 &       0 & 7{,}874{,}141 & 100\%   \\
Output tokens      &  125{,}654 &        406 &    126{,}060 &  99.7\% \\
Cost (USD)         &   \$11.481 &    \$0.052 &     \$11.533 &  99.5\% \\
API requests       &        266 &         22 &         288 &  92.4\% \\
\bottomrule
\end{tabular}
\end{table}

% ---------------------------------------------------------------------------
\subsection{Observations}\label{sec:mcp-obs}

The following observations are derived strictly from the MCP run results.

\begin{enumerate}

\item \textbf{High accuracy with zero filesystem access.}
      The model scored 163/176 (92.6\%) using only the ByteBell
      knowledge graph for codebase understanding.  13 of 22 tasks
      received a perfect 8/8, and no task scored below 5/8 ---
      eliminating the outright failure seen in the raw runs.

\item \textbf{Root cause identification matches raw-run quality.}
      Root cause scored 64/66 (97.0\%), identical to the raw runs.
      The two deductions came from tasks~15 (\texttt{14598},
      misidentifying writer vs.\ parser side) and~19 (\texttt{7606},
      incomplete --- only one of two required locations).
      The knowledge graph provides sufficient context for near-perfect
      diagnosis.

\item \textbf{File identification remains perfect.}
      All 22 tasks scored 44/44 on file identification.
      ByteBell's semantic indexing surfaces the correct file paths
      without requiring the model to manually traverse the directory
      tree.

\item \textbf{Patch quality is the bottleneck, slightly below raw.}
      The patch dimension scored 55/66 (83.3\%), marginally above
      the raw runs' 54/66 (81.8\%).
      Three tasks scored $\le$6/8: task~6 (\texttt{13579},
      over-engineered two-pass approach), task~15 (\texttt{14598},
      writer-side fix instead of parser-side), and task~19
      (\texttt{7606}, only one of two \texttt{\_\_eq\_\_} locations
      fixed).

\item \textbf{Haiku is a single-call triage layer in MCP --- even
      more extreme than raw.}
      Haiku made exactly \textbf{one} API request per task on all 22
      tasks, with \textbf{zero} cache read and \textbf{zero} cache write
      tokens.  All codebase exploration, MCP tool invocations, and patch
      reasoning were handled exclusively by Opus.  In the raw runs,
      Haiku occasionally performed multi-request exploration (up to 39
      requests on task~21); in MCP mode it is a pure routing stub.

\item \textbf{MCP is 29\% cheaper per task than raw.}
      Average cost per task is \$0.52 (MCP) versus \$0.73 (raw), a
      29\% reduction.  Total cost across 22 tasks is \$11.53 (MCP)
      versus \$16.16 (raw).  The saving comes from the knowledge graph
      enabling more targeted exploration: effective weighted input
      averages 78{,}088 tokens per task (MCP) versus 130{,}288 (raw),
      a 40\% reduction.

\item \textbf{MCP uses 1.7$\times$ fewer effective weighted input
      tokens.}
      Total effective weighted input is 1.72\,M (MCP) versus 2.87\,M
      (raw).  The knowledge graph's semantic search surfaces relevant
      code regions directly, reducing the need for broad codebase
      scanning that drives up cache reads in raw runs.

\item \textbf{Cache reads are 47\% lower in MCP.}
      Total cache reads are 7.87\,M (MCP) versus 14.79\,M (raw).
      This is the single largest contributor to MCP's cost advantage:
      fewer iterative file reads means fewer tokens billed at the
      cache-read rate.

\item \textbf{Output tokens are 55\% lower in MCP.}
      Total output is 126{,}060 (MCP) versus 276{,}939 (raw).
      MCP answers tend to be concise descriptions or targeted code
      blocks, while raw answers frequently include full unified diffs
      with tests.  The simpler answer format (no explicit diff
      instruction) and the knowledge graph's focused context both
      contribute to shorter outputs.

\item \textbf{The hardest task (13398) was solved with a comprehensive
      description.}
      Task~4 (\texttt{astropy-13398}), rated 1--4\,h human difficulty,
      required adding topocentric ITRS transforms across 6+ files
      including a new module.  The MCP answer correctly identified all
      files, described every required change (location attributes, new
      transform file, GCRS helper, CIRS/TETE propagation), and scored
      8/8 --- at \$0.70 and 281\,s, substantially cheaper and faster
      than the raw run (\$1.77, 612\,s).

\item \textbf{No outright failures.}
      Unlike the raw runs which had one fail (task~19,
      \texttt{astropy-7606}, scoring~3/8), the MCP run has no task
      below 5/8.  The same task scored 6/8 under MCP --- the knowledge
      graph surfaced enough context to fix one of the two required
      \texttt{\_\_eq\_\_} locations, whereas the raw run produced no
      patch at all.

\item \textbf{Wall-clock time is comparable to raw.}
      Average time is 249\,s (MCP) versus 251\,s (raw) --- effectively
      identical.  MCP's token savings do not translate to wall-clock
      speedup, likely because MCP tool-call round-trips offset the
      reduced token volume.  Two outliers (\texttt{14369} at 845\,s
      and \texttt{7671} at 785\,s) dominate the tail.

\end{enumerate}
